\chapter{Examples of slicing, jumps and cancellation}\label{sec:examples-current-slice-limits}
These calculations distinguish flat convergence from mass convergence and show why exceptional levels and integrality hypotheses cannot be omitted. They use the scalar slicing conventions of \cite{Sim18}, Chapter 6, Section 4.

\begin{example}[Endpoint jumps for the rectangle]\label{ex:rectangle-slice-endpoints}
Let $T=\llbracket[0,1]^2,e_1\wedge e_2\rrbracket$ and $f(x,y)=x$, as in Example \ref{ex:rectangle-slice}. Denote by $I_t$ the upward-oriented segment $\{t\}\times[0,1]$. At the two endpoints,
\[
 B_0^-=0,\qquad B_0^+=I_0,\qquad
 B_1^-=I_1,\qquad B_1^+=0.
\]
Thus the lower slice jumps from the right at zero and is left continuous at one. The two jumps come entirely from boundary mass on the cutting level, since $\|T\|(\{f=0\})=\|T\|(\{f=1\})=0$.
\end{example}
\begin{proof}
At zero the open sublevel contains no mass of $T$ or $\partial T$, so $B_0^-=0$. The closed sublevel contains no two-dimensional mass, but the boundary restriction is the left edge with downward orientation, namely $-I_0$. Hence $B_0^+=I_0$. At one, the open restriction of $T$ is all of $T$, while the boundary restriction is $\partial T-I_1$, giving $B_1^-=I_1$. The closed sublevel contains both currents completely, so $B_1^+=0$. These direct computations agree with the jump identity in Theorem \ref{thm:slice-flat-one-sided}. The interior values $B_t^-=B_t^+=I_t$, $0<t<1$, were computed in Example \ref{ex:rectangle-slice}.
\end{proof}

\begin{example}[A closed current can jump at a flat cap]\label{ex:closed-current-flat-cap-jump}
In $\mathbb R^3$, let $H$ be the outward-oriented upper unit hemisphere and let $D$ be the unit disk in $\{z=0\}$ oriented by $e_1\wedge e_2$. The integral current $T=\llbracket H\rrbracket-\llbracket D\rrbracket$ is a cycle. For $f(x,y,z)=z$, its level-zero slices satisfy
\[
 B_0^-=0,\qquad B_0^+=-\llbracket\partial D\rrbracket.
\]
Here $\partial T=0$, so the jump comes entirely from the positive mass $\|T\|(\{f=0\})=\pi$ of the flat cap. For $0<t<1$, the slice is the clockwise circle of radius $\sqrt{1-t^2}$ at height $t$, and its mass is $2\pi\sqrt{1-t^2}$.
\end{example}
\begin{proof}
Stokes' theorem gives $\partial\llbracket H\rrbracket=\partial\llbracket D\rrbracket$, with counterclockwise boundary orientation viewed from above. Thus $\partial T=0$. The strict sublevel at zero contains no current mass, while the closed sublevel restriction is $-\llbracket D\rrbracket$, giving the claimed two values. For $0<t<1$, the part above $t$ is an outward-oriented spherical cap, whose boundary is the counterclockwise horizontal circle. Since the full current is closed, the boundary of the sublevel restriction is its negative. This proves the orientation and the length formula. The flat-norm jump is nonzero because its current is nonzero, by Proposition \ref{prop:real-flat-control}; it is at most $\pi$, using the disk as a filling. The example shows why both $\|T\|$ and $\|\partial T\|$ occur in the exceptional-level measure.
\end{proof}

\begin{example}[An integral current with a dense set of slice jumps]\label{ex:dense-slice-jumps}
There exists a compactly supported integral two-current in $\mathbb R^2$, with finite mass and finite boundary mass, whose lower slices by the coordinate $f(x,y)=x$ have a dense set of discontinuities in $(0,1)$ in real flat norm. This is consistent with their one-sided continuity and bounded total flat variation.
\end{example}
\begin{proof}
Enumerate the rationals in $(0,1)$ without repetitions as $q_j$, $j\geq1$. Put $h_j=2^{-j-2}$ and
\[
 Q_j=[q_j,q_j+h_j]\times[2^{-j},2^{-j}+h_j],\qquad
 T=\sum_{j=1}^\infty\llbracket Q_j,e_1\wedge e_2\rrbracket.
\]
The vertical bands of the squares are disjoint. All squares lie in a fixed compact rectangle, and
\[
 \sum_j\mathbf M(\llbracket Q_j\rrbracket)=\sum_j h_j^2=1/48,
 \qquad
 \sum_j\mathbf M(\partial\llbracket Q_j\rrbracket)=4\sum_jh_j=1.
\]
The series and its boundary therefore converge in mass. Their explicit disjoint oriented area and edge representations have integer multiplicity, so $T$ is integral and normal. The accumulation line $y=0$ contributes no mass to either representation and no extra boundary: evaluating on a test and its derivative commutes with the absolutely convergent series.

For any $t$, $T\llcorner\{x=t\}=0$ because a vertical line has zero two-dimensional area. At $t=q_j$, the boundary restriction includes the downward left edge of $Q_j$. All other vertical edges at the same level lie in disjoint vertical bands. The horizontal edges meet that level in sets of zero length. Hence $(\partial T)\llcorner\{x=q_j\}$ is a nonzero current: a test localized near an interior point of the indicated left edge detects it and avoids all other bands. The jump formula in Theorem \ref{thm:slice-flat-one-sided} now gives a nonzero jump at every $q_j$. Nondegeneracy of the real flat norm makes each an actual flat discontinuity. Their set is dense. Proposition \ref{prop:slice-flat-total-variation} still bounds the sum of the flat jump sizes by the finite total mass and boundary mass. Thus countability of exceptional levels does not imply that they are isolated.
\end{proof}

\begin{proposition}[The exact flat norm of two opposite point masses]\label{prop:two-point-real-flat}
For $p,q\in\mathbb R^N$,
\[
 \mathcal F_{\mathbb R}(\delta_p-\delta_q)=\min\{|p-q|,2\}.
\]
The norm is the real-current flat norm of Definition \ref{def:real-current-flat-norm}; no identification with an integer-polyhedral completion is required. This is an elementary example of the flat estimates accompanying the current operations in \cite{Sim18}, Chapter 6, Section 2.
\end{proposition}
\begin{proof}
An oriented segment from $q$ to $p$ fills the difference with mass $|p-q|$. Alternatively take the difference itself as residual, with mass at most two. These decompositions prove the upper bound. If $p=q$, the conclusion follows immediately.

For the lower bound, set $e=(p-q)/|p-q|$ and $c=(p+q)/2$. The function
\[
 a(x)=\max\{-1,\min\{1,e\cdot(x-c)\}\}
\]
has absolute value at most one, Lipschitz constant at most one, and $a(p)-a(q)=\min\{|p-q|,2\}$. Mollification gives smooth functions with the same two bounds and uniform convergence to $a$. Multiply them by cutoffs equal to one on a large ball containing $p,q$, bounded by one, and with gradients tending uniformly to zero as the ball grows. The resulting compactly supported tests $a_j$ satisfy $\|a_j\|_\infty\leq1$, $\|Da_j\|_\infty\leq1+o(1)$ and the desired limiting difference of values. Proposition \ref{prop:real-flat-control} applied to $a_j/(1+o(1))$ proves the lower bound. This argument uses only the test bound for flat decompositions, not a general dual representation theorem for the flat norm.
\end{proof}

\begin{example}[A circle loses slice mass continuously in flat norm]\label{ex:circle-slice-mass-loss}
Let $T$ be the counterclockwise integration current of the circle of radius $R>0$ in $\mathbb R^2$, and set $f(x,y)=x$. For $|t|<R$, put $h_t=\sqrt{R^2-t^2}$. Then
\[
 S_t=\delta_{(t,-h_t)}-\delta_{(t,h_t)},\qquad
 \mathbf M(S_t)=2,\qquad
 \mathcal F_{\mathbb R}(S_t)=\min\{2h_t,2\}.
\]
At $t=\pm R$ both one-sided distributional slices are zero. Therefore $S_t\to0$ in flat norm as $t\to R$ from below, while the masses remain two. The coarea identity reads
\[
 \int_{\mathbb R}\mathbf M(S_t)\,dt=4R
 =\int_{\operatorname{spt}T}|\nabla^Tf|\,d\mathcal H^1.
\]
This illustrates scalar current slicing and the lower semicontinuity, rather than continuity, of mass.
\end{example}
\begin{proof}
Parametrize the circle by $\gamma(\alpha)=(R\cos\alpha,R\sin\alpha)$ with increasing $\alpha$. Its unit tangent is $(-\sin\alpha,\cos\alpha)$. The tangential derivative of $f$ is positive at the bottom point and negative at the top point, giving the stated signs by Lemma \ref{lem:slice-unit-orientation}. The points are distinct for $|t|<R$, hence their total variation is two. Their distance is $2h_t$, so Proposition \ref{prop:two-point-real-flat} gives the exact flat norm.

The boundary of the circle current vanishes. At $R$ the strict sublevel contains the whole circle except one zero-mass point, so its restriction equals $T$ and its boundary is zero. At $-R$ the strict sublevel restriction is zero. Closed sublevels give the same currents because the endpoint singletons have zero length. This proves the endpoint statements directly. Finally the mass integral is twice the interval length $2R$. In the angular parametrization the other side is $R\int_0^{2\pi}|\sin\alpha|\,d\alpha=4R$. No value assigned at the two endpoints affects this integral.
\end{proof}

\begin{example}[A level atom need not create a slice jump]\label{ex:level-atom-no-jump}
Let $T$ be the rightward integration current of $[0,1]\times\{0\}$ and let $f(x,y)=y$. Then $B_t^-=B_t^+=0$ for every $t$, although
\[
 (\|T\|+\|\partial T\|)(\{f=0\})=3.
\]
Thus the countable atomic set in Theorem \ref{thm:slice-flat-one-sided} can be strictly larger than the actual jump set.
\end{example}
\begin{proof}
The mass of $T$ is one, and its boundary consists of two opposite endpoint masses with total variation two. Both measures are supported on $f=0$. Every strict or closed sublevel either includes the entire pair $T,\partial T$ or includes neither. In the former case the slice is $\partial T-\partial T=0$; in the latter it is zero as well. Equivalently the jump formula at zero is $\partial(T\llcorner\{f=0\})-(\partial T)\llcorner\{f=0\}=0$. Tangentially, $f$ is constant on the segment and its derivative is zero.
\end{proof}

\begin{example}[Diffuse slices of a diffuse normal current]\label{ex:diffuse-current-slices}
Let $T$ be the normal two-current on $[0,1]^4$ from Example \ref{ex:diffuse-two-current}, and take $f(x)=x_1$. For $0<t<1$, its distributional slice is the one-current
\[
 B_t(\eta)=\int_{[0,1]^3}\eta_2(t,x_2,x_3,x_4)\,dx_2dx_3dx_4.
\]
Its mass is one and its boundary mass is two. It is normal but not a rectifiable one-current. Thus slicing a normal current does not in itself produce integral or rectifiable slices.
\end{example}
\begin{proof}
For $\eta=\sum_i\eta_i dx_i$, evaluate $\partial(T\llcorner\{x_1<t\})(\eta)$ by integrating
$\partial_1\eta_2-\partial_2\eta_1+\partial_3\eta_4-\partial_4\eta_3$
over $[0,t]\times[0,1]^3$. The fundamental theorem of calculus gives the new face at $x_1=t$ plus the old faces with their restricted ranges. Subtracting $(\partial T)\llcorner\{x_1<t\}$ cancels exactly those old faces. The remaining term is the stated formula.

Its polar vector is $e_2$ and its mass measure is three-dimensional Lebesgue measure on the slice cube, so the mass is one. Applying its boundary to a scalar test differentiates only in $x_2$, producing the difference of two unit two-dimensional face measures; these are disjoint, giving boundary mass two. A countably one-rectifiable subset of that three-dimensional affine hyperplane has zero three-dimensional volume, by Lemma \ref{lem:lipschitz-image-ambient-null} and the same null-set cover argument as in Proposition \ref{prop:diffuse-current-not-rectifiable}. Its mass therefore cannot be concentrated on any such set. This is an explicit normal-current calculation, and it does not apply the geometric slicing theorem outside its integer-rectifiability hypotheses.
\end{proof}
